\scanpage{26}

\deleted{Soit}
Cas d'un schéma en groupes \inserted{séparé} étale de t.f. sur $S$
dont le centre n'est pas représentable
($S={}$spectre d'un anneau de valuation
discret $A$)

On part d'un groupe ensembliste $E$,
de centre trivial, et d'un sous-groupe commutatif $\deleted{\unsure{E|H}}F$ non \deleted{\unsure{nul}} \inserted{trivial} (p.ex. le groupe engendré
par un élément $g\ne e$ de $E$). Soit
$\deleted{\unsure{F|H|\mathcal F}}G$ le sous-groupe ouvert de $E_S$ obtenu
en enlevant de la fibre spéciale
$E_k$ de $E_S$ le fermé $E_k-\deleted{\unsure{E|H}}F_k$. Donc

\[
G_0=F_k,\qquad G_1=E_K
\]

Alors le foncteur $\operatorname{Cent}(G)$ n'est pas
représentable, car il ne commute
pas à la limite projective adique
des $A_n=A/\mathfrak m^{n+1}$, car on trouve

\[
\begin{aligned}
Z(\hat A)&=\{e\}\\
\varprojlim Z(A_n)&=F
\end{aligned}
\]

\uline{N.B.} Si $G$ est \inserted{plat} de présentation finie sur $S$, à fibres connexes,
\unsure{alors} $Z$ est représentable ......

% archive sheet number: 25
